\documentclass[11pt,a4paper]{article}
\usepackage[margin=25mm]{geometry}
\usepackage{amsmath,amssymb,amsthm,hyperref}
\setlength{\emergencystretch}{2em}
\newtheorem{lemma}{Lemma}
\newtheorem{theorem}[lemma]{Theorem}
\newtheorem{corollary}[lemma]{Corollary}
\newcommand{\gen}[1]{\langle#1\rangle}
\title{Equivariant orbit aggregation for paired cuts}
\author{Human Compression project: L044 continuation}
\date{}
\begin{document}\maketitle
Let $H=\gen{s_0,\ldots,s_{r-1}}$ act on a finite set $U$, with
each $s_i$ an involution or the identity and $[s_i,s_j]=1$ for
$|i-j|>1$. Write $O_1,\ldots,O_c$ for its actual orbits.
The full local omission on $O_a$ is
\[
 H_{a,j}=\gen{s_i|_{O_a}:i\ne j}.
\]
Assign pairwise disjoint colour sets $F_1,\ldots,F_c$ to these
orbits. For every $j\in F_a$, assume that $s_j|_{O_a}$ has at
least two edges whose endpoints belong to different $H_{a,j}$-orbits.
Set $d_a=|F_a|-|O_a|/2$. The accepted selected-colour theorem,
applied to the transitive full action on each orbit, gives $d_a\le0$.
The assigned sets need not exhaust all colours unless explicitly stated.

\begin{lemma}[Transport and aggregation]\label{transport}
Suppose $\phi:O_a\longrightarrow O_b$ is a bijection satisfying
$\phi(v^{s_i})=\phi(v)^{s_i}$ for every $i$ and $v\in O_a$.
Put $m=|O_a|=|O_b|$, with $a\ne b$. Then
\[
 |F_a|+|F_b|\le\lfloor m/2\rfloor,
 \qquad d_a+d_b\le-\lceil m/2\rceil.                 \tag{1}
\]
If $m=3$, then $F_a=F_b=\varnothing$ and $d_a+d_b=-3$.
\end{lemma}
\begin{proof}
The bijection conjugates each marked restricted generator and hence
every full local omission $H_{a,j}$ to $H_{b,j}$. It carries their
orbit partitions to one another. For each colour in $F_b$, pull its
two eligible edges back to $O_a$. Together with those already chosen
for $F_a$, these give two eligible edges for every colour in
$F_a\cup F_b$ on $O_a$. Disjointness gives
$|F_a\cup F_b|=|F_a|+|F_b|$; without that hypothesis this equality
would not follow. The selected-colour theorem now bounds the union
by $\lfloor m/2\rfloor$, proving (1).
An involution on three points has at most one edge, so no colour can
be assigned to an orbit of size three under the standing hypothesis.
\end{proof}

\begin{theorem}[An exchanged odd orbit pair]\label{odd}
Suppose an involution $\delta$ centralizes $H$ and exchanges two
distinct nontrivial $H$-orbits $O_a,O_b$ of odd size. Then
\[
                         \sum_{s=1}^c d_s\le-3.       \tag{2}
\]
More explicitly, if their common size is $m\ge5$, their combined
deficit is at most $-(m+1)/2$; for $m=3$ it is $-3$.
\end{theorem}
\begin{proof}
Centralization makes $\delta|_{O_a}$ the marked equivariant
bijection of Lemma~\ref{transport}. Nontriviality and odd size
give $m\ge3$. The lemma proves the stated deficit in both ranges;
all remaining $d_s$ are nonpositive. Summing proves (2).
On the exchanged pair $\delta$ is a product of exactly $m$
transpositions, so this is precisely an odd exchanged-orbit
contribution to its sign. An exchanged pair of singleton orbits
would instead contribute only $-1$ and is not covered by the theorem.
\end{proof}

\begin{corollary}[Exact numerical tail interface]\label{tail}
Suppose the standing action is the restriction of a string tail
$\gen{\rho_{t+1},\ldots,\rho_{r-1}}$ to its support $U$.
Assume the assigned sets partition all $r-1-t$ tail indices, the
full local two-edge hypotheses hold, and Theorem~\ref{odd} applies.
Then
\[
                       r-1-t\le |U|/2-3.             \tag{3}
\]
In the two source head configurations, $t=1$ with $|U|\le n-2$,
or $t=4$ with $|U|\le n-7$, this implies $r\le(n-3)/2$.
\end{corollary}
\begin{proof}
Partitioning gives $r-1-t=\sum_s|F_s|$, so (2) gives (3).
The first configuration yields $r\le(n-4)/2$ and the second
$r\le(n-3)/2$, as required. No assertion about a particular
source assignment supplies the local eligibility hypotheses here.
\end{proof}

\paragraph{Source interface and attribution.}
Cameron--Fernandes--Leemans--Mixer, \emph{Highest rank of a polytope
for $A_n$}, arXiv:1605.09173v1, physical page 42, Proposition 5.31,
case (1), uses an odd pair of exchanged nontrivial tail orbits.
The orbit assignment and deficit framework are borrowed from its
Propositions 5.27 and 5.31. The present argument replaces the
maximal-empty-assignment step by a direct marked-equivariance
aggregation, with a bound valid for every odd orbit size.
The full accepted paired-cut and equality proofs remain paid through
the selected-colour theorem; this is not an independent reproof of
those results. No projected intersection property is used.

\paragraph{Scope.}
This proves the algebraic deficit route in case (1), conditional on
the exact local hypotheses above. Full local omissions are the
images of the tail omissions on each orbit; they are not silently
identified with omissions of an ambient larger group or of the
selected-colour subgroup. No assertion is made that the source's
other orbit configurations meet the two-edge hypothesis. The
internal-orbit cases (2)--(4), their equality rigidity and the
remaining source tables are not resolved here. There is no new
finite computation, formal theorem, whole-original recovery,
conjecture solution or priority claim.
\end{document}
